Kutatásunk során többek között rátaláltunk a Budapest Go-ra, talán ez a
meglévő megoldás áll a legközelebb a mi ötletünkhöz. Céljaink hasonlóak: hogy
a rendszer egyetlen platformon biztosítsa a tömegközlekedéshez kapcsolódó
különféle szolgáltatásokat, beleértve a jegyvásárlást, útvonaltervezést és a valós
idejű járatinformációk megjelenítését. Lényeges különbség viszont az, hogy a
Budapest Go nem rendelkezik adminisztrációs platformmal, amely lehetővé
tenné a járatok, járműflotta, sofőrök és útvonalak közvetlen menedzselését,
vagy legalábbis sehol nem jelenik meg róla információ.

Emellett a Budapest Go jelenleg nem kínál olyan POS terminál alkalmazást,
amelyet közvetlenül a buszokon vagy villamosokon használnának jegyek és
bérletek fizetésére. Az alkalmazás NFC-alapú mobilfizetést kínál, de a fizikai POS
terminál integráció nem része a rendszernek. Bár az utasok számára valós idejű
járatkövetés elérhető, a Budapest Go nem kínál olyan funkciót, amely a
járművezetők vagy a közlekedési vállalatok számára biztosítana egy külön
alkalmazást a járművek helyzetének követésére és kezelésére.

A Budapest Go egy általános közlekedési rendszer, amely a városi közlekedés
minden elemét lefedi. Nincs kifejezetten busztársaságokra szabott adminisztrációs felülete, amely például specifikus buszútvonalakat, menetrendeket vagy buszvezetői beosztásokat kezelne. Mi viszont pont erre fektetnénk a hangsúlyt és egy olyan személyre szabott rendszert hoznánk létre, amelyet bármely busztársaság tudna implementálni.